% SOURCE_AUTHORITY: sga5_fr_workpass.tex, lines 7737--7774 (LF-normalized unit).
% SOURCE_SHA256: 791F4EFFC5E02832D5D77ED03518C8156D6F07E4C8238B03545DB93D883FBB28
\begin{statement}{Proposição 5.2.2}
Os objetos da categoria abeliana AR-\(J\)-adn\((\mathcal C)\) são noetherianos.
\end{statement}

\begin{prooftext}
Reduzimo-nos imediatamente a verificar que, se se têm um objeto \(J\)-ádico noetheriano \(X\), uma sucessão de objetos \(J\)-ádicos noetherianos \((Y^n)_{n\in\mathbb N}\) e morfismos de \(\Hom(\mathbb N^\circ,\mathcal C)\)
\[
v_n:Y^n\longrightarrow X
\]
tais que
\begin{enumerate}[label=\textup{(\roman*)}]
\item Os \(v_n\) se tornam monomorfismos em AR-\(J\)-adn\((\mathcal C)\);
\item para todo inteiro \(n\), \(v_n\) se fatora através de \(v_{n+1}\) em AR-\(J\)-adn\((\mathcal C)\),
\end{enumerate}
então os \(v_n\) se estabilizam, salvo isomorfismo, para \(n\) suficientemente grande, em AR-\(J\)-adn\((\mathcal C)\).

Ora, dizer que \(v_n\) se fatora através de \(v_{n+1}\) em AR-\(J\)-adn\((\mathcal C)\) significa que existem um inteiro \(r\) e um morfismo
\[
f:Y^n[r]\longrightarrow Y^{n+1}
\]
que, se \(a\) designa o morfismo canônico
\[
Y^n[r]\longrightarrow Y^n,
\]
verifica-se a igualdade que se segue:
\[
v_n\circ a=v_{n+1}\circ f .
\]
Como o morfismo \(a\) é estrito, deduz-se em particular que a imagem de \(v_n\) está contida na de \(v_{n+1}\). Substituindo cada \(Y^n\) por \(\Im(v_n)\), reduzimo-nos a verificar que toda sucessão crescente de subsistemas projetivos estritos de \(X\) é estacionária. Ora, se \((Z^n)_{n\in\mathbb N}\) é tal sucessão, como \(\grs(X)\) é noetheriano, a sucessão crescente dos \(\grs(Z^n)\) é estacionária e, portanto, também o é a dos \(Z^n\).
\end{prooftext}

Resumamos:

\begin{statement}{Teorema 5.2.3}
As categorias \(J\)-adn\((\mathcal C)\) e AR-\(J\)-adn\((\mathcal C)\) são abelianas e noetherianas.
\end{statement}

Quanto à AR-\(J\)-adn\((\mathcal C)\), é o que acabamos de ver. Quanto a \(J\)-adn\((\mathcal C)\), é equivalente a essa.
